\section{Experiments}
\label{sec:experimental_setup}

\subsection{Data Splits}
\label{sec:dataset_splits}

We evaluate YogaLLaVA V3 on the combined domain-knowledge and
geometry-grounded QA corpus introduced in Sec.~\ref{sec:dataset}. The fixed
splits contain 10,090 records and 10,177 supervised angle slots
(Table~\ref{tab:dataset_splits}). The geometry-grounded subset is constructed
from 3,000 images, with two QA records per image, and comprises 3,000
direct-angle, 905 comparison, 900 regional-summary, 600 misconception, and 595
instructor-observation examples. Test images are excluded from training, while
pose categories are shared across splits. The resulting protocol evaluates
generalization to unseen images of known asanas rather than to unseen pose
categories.

\begin{table}[t]
    \centering
    \caption{Composition of the instruction-tuning corpus.}
    \label{tab:dataset_splits}
    \resizebox{\columnwidth}{!}{
    \begin{tabular}{lrrrr}
        \toprule
        \textbf{Split} &
        \textbf{Geometry QA} &
        \textbf{Domain QA} &
        \textbf{Total} &
        \textbf{Angle Slots} \\
        \midrule
        Train      & 5,396 & 3,690 & 9,086 & 9,163 \\
        Validation &   294 &   200 &   494 &   493 \\
        Test       &   310 &   200 &   510 &   521 \\
        \midrule
        Total      & 6,000 & 4,090 & 10,090 & 10,177 \\
        \bottomrule
    \end{tabular}}
\end{table}

\subsection{Implementation Details}
\label{sec:implementation_details}

YogaLLaVA V3 is initialized from LLaVA-1.5-7B. Its CLIP ViT-L/14-336
image encoder operates at $336\!\times\!336$ resolution, and the multimodal
projector maps the visual features to the 4,096-dimensional language space.
The maximum sequence length is 1,280 tokens. Both the vision encoder and
multimodal projector remain frozen; the language decoder is adapted using the
LoRA settings described in Sec.~\ref{sec:training_objective}, together
with the proposed geometry and angle modules.

We train for 12 epochs using AdamW with weight decay 0.01 and gradient clipping
at 1.0. A per-device batch size of four and eight-step gradient accumulation
produce an effective batch size of 32. We use a cosine learning-rate schedule
with linear warm-up over the first 3\% of optimizer updates. The learning rates
are $1.5\!\times\!10^{-4}$ for LoRA parameters,
$2\!\times\!10^{-5}$ for the trainable input and output embeddings, and
$1\!\times\!10^{-4}$ for the geometry encoder, angle-feedback encoder, angle
embedding, and regression head. The pretrained backbone is loaded in bfloat16,
whereas the added numeric modules are kept in float32. Training is
performed on one NVIDIA RTX A6000 GPU with fixed Python, NumPy, PyTorch, and
CUDA seeds.

We evaluate the model after every epoch and select the checkpoint with the
lowest teacher-forced validation angle RMSE. The selected checkpoint obtains a
validation RMSE of $19.83^\circ$, MAE of $12.63^\circ$, and $R^2$ of 0.874.
We report only the selected checkpoint because intermediate epoch-wise values
are optimization diagnostics rather than comparative experimental results.

\subsection{Evaluation Protocol}
\label{sec:evaluation_protocol}

\noindent\textbf{Teacher-forced angle regression.}
Our primary numeric evaluation supplies the reference answer structure,
including the number, positions, and ordering of \texttt{<ANGLE>} tokens, and
evaluates the continuous prediction at each known slot. Given $N$ reference
angles $\theta_i$ and predictions $\widehat{\theta}_i$, we report MAE, RMSE,
$R^2$, Pearson correlation, and the fractions of predictions within
$5^\circ$, $10^\circ$, and $15^\circ$ of the references:
\begin{equation}
\operatorname{MAE}=\frac{1}{N}\sum_{i=1}^{N}
\left|\widehat{\theta}_i-\theta_i\right|,
\qquad
\operatorname{RMSE}=\sqrt{\frac{1}{N}\sum_{i=1}^{N}
(\widehat{\theta}_i-\theta_i)^2}.
\label{eq:evaluation_errors}
\end{equation}
Teacher forcing isolates continuous prediction given the correct linguistic
slot structure; it does not evaluate whether the model can autonomously emit,
order, and contextualize the required slots.

\noindent\textbf{Autoregressive generation.}
We therefore additionally evaluate free-running generation on all 310
geometry-grounded test records. YogaLLaVA V3 receives the image, question, and
structured geometry. LLaVA-1.5-7B and Qwen2.5-VL-7B-Instruct~\cite{bai2025qwen25vl} are evaluated
zero-shot using only the image and question; the questions specify the target
metrics and their expected order but do not reveal numeric values. We report
(i) slot-count accuracy over the complete test set and (ii) numeric error on
the common-success subset for which all models generate the correct number of
values. This distinction avoids directly comparing numeric errors computed on
different model-specific subsets.

All numeric references are pseudo-measurements derived from SAM 3D Body rather
than motion capture or manually verified biomechanical annotations. Moreover,
518 of the 521 teacher-forced test targets are present in the structured
geometry input. The primary evaluation therefore measures question-conditioned
descriptor selection, slot filling, reproduction, and linguistic expression, not
independent physical angle recovery from RGB images.
